%!TEX root = ../main.tex
\section{Conclusion}
\label{sec:conclusion}

Embodied intelligence does not learn from raw episode counts; it learns from the
distribution produced after a sequence of often hidden decisions.  By organizing
those decisions along operation, signal, utility evidence, and lifecycle axes,
this survey connects acquisition, alignment, semantic enrichment, curation,
mixture design, synthesis, validation, and deployment feedback within one
problem.  The synthesis rejects two convenient simplifications: data quality is
not a policy-independent scalar, and generated plausibility is not physical or
training validity.

The most consequential shift is from one-shot cleaning to target-conditional
feedback control.  This shift makes stronger evidence possible---the policy can
reveal which data it needs---but also creates error-amplification, collapse, and
safety risks.  Progress therefore depends on calibrated links between cheap
proxies and controlled policy utility; transition-level, cross-embodiment
semantics; independent layered verification; structured failure and recovery;
and joint accounting of preparation, training, and robot cost.  A shared raw-pool
benchmark and versioned lineage would turn these principles into comparable
experiments.  Under that discipline, data preparation can become an independent,
reproducible research layer for embodied foundation models rather than an
unreported collection of engineering choices.
